\section{Limitations}
\label{sec:limitations}

\begin{enumerate}[leftmargin=*, label=\textbf{L\arabic*.}]

\item \textbf{Scale.}
The tested KGs contain 237--536 nodes.
Production biomedical KGs (e.g., PrimeKG~\cite{TODO-PrimeKG},
Hetionet~\cite{TODO-Hetionet}) contain millions of nodes and edges.
Behaviour of the pipeline at production scale is unknown; in particular,
the filter's performance and the top-20 composition may change significantly
with increased graph density.

\item \textbf{Semi-manual bridge alignment.}
Bridge entities were selected by domain knowledge, not discovered algorithmically.
Automated bridge discovery (e.g., via entity linking to a shared ontology) would
remove researcher degrees of freedom from bridge selection and would be required
for a production deployment.

\item \textbf{Single-reviewer qualitative labels.}
Run~011 labels (the sole source of ground truth for Claim~2's before-filter rates)
were assigned by one reviewer without inter-rater reliability measurement.
Future work should obtain independent labels from domain experts in each domain pair.

\item \textbf{Post-hoc filter derivation.}
The drift filter was derived from Run~011 output.
Run~013 provides generalisation evidence, but a prospective filter design
(derived from a held-out dataset) would provide stronger overfitting protection.

\item \textbf{No baseline comparison.}
No competing cross-domain discovery system was evaluated.
The paper cannot claim superiority to existing approaches; it characterises
the behaviour of the proposed mechanism.

\item \textbf{Three domain pairs only.}
The evidence for Claim~3 (bridge dispersion predicts yield) rests on 3 data points.
Statistical characterisation of the bridge dispersion effect requires testing
across a larger diversity of domain pairs.

\item \textbf{Top-20 depth problem (unsolved).}
In all configurations, top-20 scored candidates are exclusively 2-hop paths.
Deep cross-domain candidates---the primary scientifically interesting output---do not
surface in practical rankings.
Depth-promotion methods, minimum-depth tiers, or separate ranking tracks would be
required for deep candidates to be actionable.

\item \textbf{H2 (noise robustness) not retested on real data.}
Noise robustness (tolerance to random edge deletion) was validated on synthetic
KGs only (Runs~001--006) and was not retested under real-data conditions.
Its applicability to noisy production KGs is unknown.

\item \textbf{Relation semantic insufficiency (central open problem).}
The class-level encoding of \texttt{produces} and \texttt{is\_substrate\_of}
relations means that filter-passing candidates can carry chemically or
biologically invalid edges that the filter cannot detect.
Resolving this requires adding substrate-specificity metadata (a
\texttt{produces\_class} distinction) and biological compartment annotations
to the KG schema.
Until such metadata is available, all candidate paths should be treated as
\textit{structurally generated cross-domain connections} requiring expert
validation before experimental follow-up.

\end{enumerate}
